\documentclass[10pt,a4paper]{amsart}
\usepackage{amsmath,amssymb,amsthm,graphicx}
% For commutative diagrams we use Xy-pic
\usepackage[all]{xy}

% We use 2cell for 2-commutative diagrams.
\xyoption{2cell}
% Theorem environments.
%
\theoremstyle{plain}
\newtheorem{theorem}[subsection]{Theorem}
\newtheorem{proposition}[subsection]{Proposition}
\newtheorem{lemma}[subsection]{Lemma}

\theoremstyle{definition}
\newtheorem{definition}[subsection]{Definition}
\newtheorem{example}[subsection]{Example}
\newtheorem{exercise}[subsection]{Exercise}
\newtheorem{situation}[subsection]{Situation}

\theoremstyle{remark}
\newtheorem{remark}[subsection]{Remark}
\newtheorem{remarks}[subsection]{Remarks}

\numberwithin{equation}{subsection}

% Macros
%
\def\lim{\mathop{\mathrm{lim}}\nolimits}
\def\colim{\mathop{\mathrm{colim}}\nolimits}
\def\Spec{\mathop{\mathrm{Spec}}}
\def\Hom{\mathop{\mathrm{Hom}}\nolimits}
\def\Ext{\mathop{\mathrm{Ext}}\nolimits}
\def\SheafHom{\mathop{\mathcal{H}\!\mathit{om}}\nolimits}
\def\SheafExt{\mathop{\mathcal{E}\!\mathit{xt}}\nolimits}
\def\Sch{\mathit{Sch}}
\def\Mor{\mathop{\mathrm{Mor}}\nolimits}
\def\Ob{\mathop{\mathrm{Ob}}\nolimits}
\def\Sh{\mathop{\mathit{Sh}}\nolimits}
\def\NL{\mathop{N\!L}\nolimits}
\def\CH{\mathop{\mathrm{CH}}\nolimits}
\def\proetale{{pro\text{-}\acute{e}tale}}
\def\etale{{\acute{e}tale}}
\def\QCoh{\mathit{QCoh}}
\def\Ker{\mathop{\mathrm{Ker}}}
\def\Im{\mathop{\mathrm{Im}}}
\def\Coker{\mathop{\mathrm{Coker}}}
\def\Coim{\mathop{\mathrm{Coim}}}

% Boxtimes
%
\DeclareMathSymbol{\boxtimes}{\mathbin}{AMSa}{"02}

%
% Macros for moduli stacks/spaces
%
\def\QCohstack{\mathcal{QC}\!\mathit{oh}}
\def\Cohstack{\mathcal{C}\!\mathit{oh}}
\def\Spacesstack{\mathcal{S}\!\mathit{paces}}
\def\Quotfunctor{\mathrm{Quot}}
\def\Hilbfunctor{\mathrm{Hilb}}
\def\Curvesstack{\mathcal{C}\!\mathit{urves}}
\def\Polarizedstack{\mathcal{P}\!\mathit{olarized}}
\def\Complexesstack{\mathcal{C}\!\mathit{omplexes}}
% \Pic is the operator that assigns to X its picard group, usage \Pic(X)
% \Picardstack_{X/B} denotes the Picard stack of X over B
% \Picardfunctor_{X/B} denotes the Picard functor of X over B
\def\Pic{\mathop{\mathrm{Pic}}\nolimits}
\def\Picardstack{\mathcal{P}\!\mathit{ic}}
\def\Picardfunctor{\mathrm{Pic}}
\def\Deformationcategory{\mathcal{D}\!\mathit{ef}}
\usepackage{cleveref}
\crefname{originalreference}{original reference}{original references}
\Crefname{originalreference}{Original reference}{Original references}
\makeatletter
\newlabel{lemma-torsion-cohomology}{{0DDC}{source}{}{sourcepointer-1}{}}
\newlabel{lemma-torsion-cohomology@cref}{{[originalreference][0][]0DDC}{[1][0][]source}{}{}{}}
\makeatother
\crefname{originallemma}{original lemma}{original lemmas}
\Crefname{originallemma}{Original lemma}{Original lemmas}
\makeatletter
\newlabel{sites-cohomology-lemma-higher-direct-images}{{072W}{source}{}{sourcepointer-2}{}}
\newlabel{sites-cohomology-lemma-higher-direct-images@cref}{{[originallemma][0][]072W}{[1][0][]source}{}{}{}}
\makeatother
\crefname{originallemma}{original lemma}{original lemmas}
\Crefname{originallemma}{Original lemma}{Original lemmas}
\makeatletter
\newlabel{sites-cohomology-lemma-unbounded-describe-higher-direct-images}{{0D6I}{source}{}{sourcepointer-3}{}}
\newlabel{sites-cohomology-lemma-unbounded-describe-higher-direct-images@cref}{{[originallemma][0][]0D6I}{[1][0][]source}{}{}{}}
\makeatother
\title{Higher direct images of a torsion sheaf, and of a complex with torsion cohomology, along a quasi-compact quasi-separated morphism are torsion}
\author{The Stacks Project Authors}
\date{stacks-project:etale-cohomology.tex:local-snapshot-20261011; GNU FDL 1.2 or later}
\begin{document}
\maketitle
\noindent\textit{Editorial note.} All delivered mathematical text is the author's own, copied from the pinned Stacks Project chapter file etale-cohomology.tex; no statement or proof step was rewritten, reordered or removed.\par\medskip
\noindent\textit{Editorial note.} The counted proof is the author's complete argument: the proof is written out in full for both parts: the presheaf description of the higher direct images, the reduction to affine V where the fibre product is quasi-compact and quasi-separated, and the appeal to the torsion statement.\par\medskip
\noindent\textit{Editorial note.} The label lemma-torsion-direct-image names the counted statement, and the proof environment immediately following it in the source is delivered verbatim as its proof.\par\medskip
\noindent\textit{Editorial note.} The counted statement is carried by the author's own lemma environment in the chapter source.\par\medskip
\noindent\textit{Editorial note.} The statement and proof are copied from the pinned chapter file; the cited inputs are kept as external pointers with their canonical Stacks tags.\par\medskip
\noindent\textit{Editorial note.} The preamble is the chapter's own native preamble excerpt already used for this source: the Xy-pic setup (native preamble lines 14-18) and the native theorem declarations with the mathematical macros (native preamble lines 154-217).\par\medskip
\noindent\textit{Allowed standard prerequisites (editorial).}\begin{itemize}
\item Description of the higher direct images by the cohomology of the base change. Cohomology on Sites, the lemmas identifying R\textasciicircum{}n f\_* of a sheaf and of a complex with the sheaf associated to V |-> H\textasciicircum{}n(X x\_Y V, -). Supplies the presheaf description used to reduce both assertions to a computation on X x\_Y V.
\item Torsion cohomology of a quasi-compact quasi-separated scheme. Same chapter, the lemma stating that the etale cohomology of a torsion sheaf on a quasi-compact quasi-separated scheme is torsion. Supplies the torsion property of the cohomology groups appearing in the presheaf description.
\end{itemize}\medskip
\ifdefined\hypertarget\hypertarget{sourcepointer-1}{}\fi
\noindent\textit{Source reference.} Original reference 0DDC: sources/stacks-etale-cohomology/pinned-original-source.tex, source lines 13899--13899. Prerequisite pointer invoked at proof line(s) 13943; it names a standard fact of the same chapter that the counted argument uses, and every step of the counted proof built on it is retained verbatim in the delivered body.\par\medskip
\ifdefined\hypertarget\hypertarget{sourcepointer-2}{}\fi
\noindent\textit{Source reference.} Original lemma 072W: sources/stacks-chow/tags.txt, source lines 9040--9040. Prerequisite pointer invoked at proof line(s) 13940; it names a standard fact of a sibling chapter that the counted argument uses, and every step of the counted proof built on it is retained verbatim in the delivered body.\par\medskip
\ifdefined\hypertarget\hypertarget{sourcepointer-3}{}\fi
\noindent\textit{Source reference.} Original lemma 0D6I: sources/stacks-chow/tags.txt, source lines 16523--16523. Prerequisite pointer invoked at proof line(s) 13950; it names a standard fact of a sibling chapter that the counted argument uses, and every step of the counted proof built on it is retained verbatim in the delivered body.\par\medskip
\begin{lemma}
\label{lemma-torsion-direct-image}
Let $f : X \to Y$ be a quasi-compact and quasi-separated
morphism of schemes.
\begin{enumerate}
\item If $\mathcal{F}$ is a torsion abelian sheaf on $X_\etale$, then
$R^nf_*\mathcal{F}$ is a torsion abelian sheaf on $Y_\etale$ for all $n$.
\item If $K$ in $D^+(X_\etale)$ has torsion cohomology sheaves, then
$Rf_*K$ is an object of $D^+(Y_\etale)$ whose cohomology sheaves are
torsion abelian sheaves.
\end{enumerate}
\end{lemma}

\begin{proof}
Proof of (1). Recall that $R^nf_*\mathcal{F}$ is the sheaf associated
to the presheaf $V \mapsto H^n_\etale(X \times_Y V, \mathcal{F})$
on $Y_\etale$. See Cohomology on Sites,
Lemma \ref{sites-cohomology-lemma-higher-direct-images}.
If we choose $V$ affine, then $X \times_Y V$ is
quasi-compact and quasi-separated because $f$ is, hence
we can apply Lemma \ref{lemma-torsion-cohomology} to see that
$H^n_\etale(X \times_Y V, \mathcal{F})$ is torsion.

\medskip\noindent
Proof of (2). Recall that $R^nf_*K$ is the sheaf associated
to the presheaf $V \mapsto H^n_\etale(X \times_Y V, K)$
on $Y_\etale$. See Cohomology on Sites,
Lemma \ref{sites-cohomology-lemma-unbounded-describe-higher-direct-images}.
If we choose $V$ affine, then $X \times_Y V$ is
quasi-compact and quasi-separated because $f$ is, hence
we can apply Lemma \ref{lemma-torsion-cohomology} to see that
$H^n_\etale(X \times_Y V, K)$ is torsion.
\end{proof}

\end{document}
